\documentclass{article}

\usepackage{amsmath,amssymb}
\usepackage{xurl}
\usepackage[hidelinks]{hyperref}

% Shared commands for notes in docs/paper-gaps/.

\DeclareUnicodeCharacter{2080}{\ifmmode{}_{0}\else\textsubscript{0}\fi}
\DeclareUnicodeCharacter{2081}{\ifmmode{}_{1}\else\textsubscript{1}\fi}
\DeclareUnicodeCharacter{2082}{\ifmmode{}_{2}\else\textsubscript{2}\fi}
\DeclareUnicodeCharacter{2099}{\ifmmode{}_{n}\else\textsubscript{n}\fi}

\newcommand{\Lean}{\textsc{Lean}}
\ExplSyntaxOn
\NewDocumentCommand{\leanid}{m}{
  \ifmmode
    \textsf{\detokenize{#1}}
  \else
    \group_begin:
    \urlstyle{sf}
    \tl_set:Nn \l_tmpa_tl {#1}
    \tl_replace_all:Nnn \l_tmpa_tl {\_} {_}
    \exp_args:NV \path \l_tmpa_tl
    \group_end:
  \fi
}
\ExplSyntaxOff
\providecommand{\tr}{\operatorname{tr}}
\newcommand{\tnleanIssueBase}{https://github.com/LionSR/TNLean/issues/}
\newcommand{\tnleanPrBase}{https://github.com/LionSR/TNLean/pull/}
\newcommand{\tnleanDocBase}{https://sirui-lu.com/TNLean/docs/}
\newcommand{\ghissue}[1]{\href{\tnleanIssueBase#1}{GitHub issue \##1}}
\newcommand{\ghpr}[1]{\href{\tnleanPrBase#1}{PR \##1}}
\newcommand{\doclink}[1]{\href{\tnleanDocBase#1.html}{\path{#1}}}

% Dirac notation and the identity operator, as in blueprint/src/macros/common.tex.
% \providecommand keeps note-local definitions authoritative.
\usepackage{dsfont}
\providecommand{\ket}[1]{|#1\rangle}
\providecommand{\bra}[1]{\langle #1|}
\providecommand{\braket}[2]{\langle #1 | #2 \rangle}
\providecommand{\ketbra}[2]{|#1\rangle\!\langle #2|}
\providecommand{\Id}{\mathds{1}}
\providecommand{\one}{\mathds{1}}
\providecommand{\C}{\mathbb{C}}
\providecommand{\MN}[1]{M_{#1}(\C)}
\providecommand{\MD}{M_D(\C)}

% External referencing for paper sources.  The build helper writes sanitized
% auxiliary files under build/paper-aux/ and excludes duplicated source labels.
\usepackage{xr}

% Import a generated paper auxiliary file with a caller-supplied label prefix.
% The candidate paths support builds from docs/paper-gaps/, the repository root,
% and build/paper-aux/ itself.
\newcommand{\importpaperaux}[2]{%
  \IfFileExists{../../build/paper-aux/#2.aux}{%
    \externaldocument[#1]{../../build/paper-aux/#2}%
  }{%
    \IfFileExists{build/paper-aux/#2.aux}{%
      \externaldocument[#1]{build/paper-aux/#2}%
    }{%
      \IfFileExists{#2.aux}{%
        \externaldocument[#1]{#2}%
      }{}%
    }%
  }%
}

% General external reference macros
\newcommand{\paperref}[2]{\ref{#1-#2}}
\newcommand{\papereqref}[2]{(\ref{#1-#2})}

% CPSV16 source (1606.00608 / MPDO-22-12-17-2.tex) external document and shortcuts
\importpaperaux{cpsv16-}{1606.00608/MPDO-22-12-17-2-tnlean}
\newcommand{\cpsvref}[1]{\paperref{cpsv16}{#1}}
\newcommand{\cpsveqref}[1]{\papereqref{cpsv16}{#1}}

% Verdict marker: \gapnote{<kind>}{<status>}. Kind is the policy
% classification (clarification, local-correction, scope-restriction,
% unfaithful, false-source, open-gap); status is open, wip (elimination
% underway), resolved, or historical. Machine-read by the site index;
% prints a verdict line.
\newcommand{\gapnote}[2]{\par\noindent\textbf{Verdict:} #1 (#2).\par}


\title{Nondegeneracy in the Alternating-Network Proportionality Lemma}
\author{}
\date{2026-08-30}

\begin{document}
\maketitle
\gapnote{local-correction}{resolved}

\section{Local nondegeneracy convention}

Lemma~\texttt{lemma:1} of
Ref.~\cite[lines~2296--2395]{FrancoRubio2025MPUGauging} considers a connected
open chain
\begin{align}
    A_1B_1A_2B_2\cdots A_nB_n
        &=A_1B'_1A_2B'_2\cdots A_nB'_n,
    \label{eq:fbc25_alternating_network_nondegeneracy_source_equality}
\end{align}
where every $A_i$ is injective.  It concludes that there are scalars
$\beta_i$ such that
\begin{align}
    B_i&=\beta_iB'_i,
    \label{eq:fbc25_alternating_network_nondegeneracy_source_proportionality}\\
    \prod_{i=1}^n\beta_i&=1.
    \label{eq:fbc25_alternating_network_nondegeneracy_source_product}
\end{align}
The formal treatment adopts the intended nondegenerate convention that every
intervening tensor $B_i$ is nonzero.  This convention is present at both
applications of the lemma in the article.

\section{Nondegenerate correction}

The corrected statement assumes
\begin{align}
    B_i&\ne0\qquad (1\leq i\leq n).
    \label{eq:fbc25_alternating_network_nondegeneracy_nonzero_factors}
\end{align}
Contracting the physical leg of each $A_i$ against an inverse tensor, exactly
as in the source proof, separates the intervening tensors and gives
\begin{align}
    \bigotimes_{i=1}^n B_i
        &=\bigotimes_{i=1}^n B'_i.
    \label{eq:fbc25_alternating_network_nondegeneracy_pure_tensor_equality}
\end{align}
Condition
(\ref{eq:fbc25_alternating_network_nondegeneracy_nonzero_factors}) makes the
left-hand side of
(\ref{eq:fbc25_alternating_network_nondegeneracy_pure_tensor_equality})
nonzero.  Proportionality of the factors of two equal nonzero pure tensors
then gives
(\ref{eq:fbc25_alternating_network_nondegeneracy_source_proportionality})--%
(\ref{eq:fbc25_alternating_network_nondegeneracy_source_product}).  A
one-sided nonvanishing assumption is sufficient: the pure-tensor equality
forces every $B'_i$ to be nonzero as well.

The two uses in Ref.~\cite{FrancoRubio2025MPUGauging} already lie in this
nondegenerate regime.  The truncated symmetry $U^L_g$ is unitary by the
statement at line~2099.  For the application at line~2567, lines~2444--2466
identify one side of the top boundary network with
$(U^L_{g^{-1}})^\dagger$.  This network is nonzero; hence none of its
constituent intervening factors can vanish.  One may take these factors as the
single-sided family $B_i$ in
(\ref{eq:fbc25_alternating_network_nondegeneracy_nonzero_factors}).  For the
application at line~3157, lines~3033--3053 identify the corresponding network
with $\sigma_g(U^L_{g^{-1}})^\dagger$, where
$\sigma_g\in\{1,-1\}$.  It is again nonzero for the same reason.  Neither use
requires a stronger invertibility assumption on every $B_i$.

\section{Formal treatment}

The formal theorem models the connected open chain with its horizontal ends
uncontracted.  Matrix units supplied to the one-sided inverses of the
injective $A_i$ retain both virtual indices, so their contraction proves
(\ref{eq:fbc25_alternating_network_nondegeneracy_pure_tensor_equality}) rather
than merely an equality of matrix products.  The final factorwise conclusion
uses the existing pure-tensor proportionality theorem.\footnote{Formal
declaration:
\leanid{MPSTensor.exists\_scalars\_of\_openAlternatingNetwork\_eq} in
\doclink{TNLean/MPS/Chain/AlternatingNetwork}.}

\textbf{Verdict: resolved local fix.}
The intended nondegeneracy convention is present at both applications in the
article.  The unrestricted printed wording is not used as a formal theorem.

\bibliographystyle{alpha}
\bibliography{references}

\end{document}
